% ===== PRÉAMBULE CANONIQUE — à coller VERBATIM en tête de CHAQUE .tex =====
\documentclass[11pt,a4paper]{article}
\usepackage[utf8]{inputenc}\usepackage[T1]{fontenc}\usepackage{lmodern}
\usepackage[french]{babel}
\usepackage{amsmath,amssymb,mathtools}\usepackage{icomma}
\usepackage[version=4]{mhchem}          % \ce{...} : équations & formules
\usepackage{siunitx}\sisetup{locale=FR} % \qty{}{}, \si{}, \ang{}
\usepackage{chemfig}                    % \chemfig{...} : structures développées
\usepackage{xcolor}\usepackage{colortbl}
\usepackage{tikz}\usepackage{pgfplots}\pgfplotsset{compat=1.18}
\usepackage[most]{tcolorbox}\usepackage{enumitem}\usepackage{array,booktabs}
\usepackage[a4paper,margin=2.2cm]{geometry}
\usetikzlibrary{arrows.meta,calc,angles,quotes,positioning,patterns,backgrounds,shapes.geometric}
\definecolor{primary}{RGB}{222,49,99}\definecolor{primarydark}{RGB}{150,22,55}
\definecolor{defcol}{RGB}{0,114,178}\definecolor{methcol}{RGB}{52,92,148}
\definecolor{excol}{RGB}{0,135,90}\definecolor{attncol}{RGB}{200,40,55}
\tcbset{boxbase/.style={enhanced,breakable,boxrule=0.9pt,arc=2.5pt,
  left=3mm,right=3mm,top=2.3mm,bottom=2.3mm,fonttitle=\bfseries,coltitle=white}}
% --- boîtes TUTORIEL (noms EXACTS, ne pas renommer) ---
\newtcolorbox{cle}[1][]{boxbase,colback=defcol!6,colframe=defcol,
  title={\textsc{À retenir}\ifx\relax#1\relax\relax\else\textmd{ : \itshape #1}\fi}}
\newtcolorbox{methode}[1][]{boxbase,colback=methcol!7,colframe=methcol,sharp corners,
  title={\textsc{Méthode}\ifx\relax#1\relax\relax\else\textmd{ --- \itshape #1}\fi}}
\newtcolorbox{exemplebox}[1][]{boxbase,colback=excol!5,colframe=excol,
  title={\textsc{Exemple}\ifx\relax#1\relax\relax\else\textmd{ : \itshape #1}\fi}}
\newtcolorbox{piege}[1][]{boxbase,colback=attncol!6,colframe=attncol,
  title={\textsc{Piège}\ifx\relax#1\relax\relax\else\textmd{ : \itshape #1}\fi}}
\newtcolorbox{remarque}[1][]{boxbase,colback=black!4,colframe=black!45,
  title={\textsc{Remarque}\ifx\relax#1\relax\relax\else\textmd{ : \itshape #1}\fi}}
% --- boîtes EXERCICE / CORRIGÉ (noms EXACTS, auto-comptées) ---
\newtcolorbox[auto counter]{exo}[1][]{boxbase,colback=primary!4,colframe=primary,
  title={\textsc{Exercice}~\thetcbcounter\ifx\relax#1\relax\relax\else\textmd{ --- \itshape #1}\fi}}
\newtcolorbox[auto counter]{corrige}[1][]{boxbase,colback=excol!4,colframe=excol,
  title={\textsc{Corrigé}~\thetcbcounter\ifx\relax#1\relax\relax\else\textmd{ --- \itshape #1}\fi}}
\newcommand{\R}{\mathbb{R}}\newcommand{\N}{\mathbb{N}}\newcommand{\Z}{\mathbb{Z}}\newcommand{\ds}{\displaystyle}
% (un \usepackage{fancyhdr}+en-tête est possible ; il est ignoré côté web)
% =========================================================================
\begin{document}

\begin{corrige}[pression totale = somme des partielles]
On additionne (ou on soustrait) les pressions partielles.
\begin{enumerate}[label=\alph*)]
  \item $p_{\text{tot}}=0{,}3+0{,}5=\qty{0.8}{atm}$.
  \item $p_B=p_{\text{tot}}-p_A=2{,}0-1{,}2=\qty{0.8}{atm}$.
  \item $p_{\text{tot}}=0{,}2+0{,}3+0{,}5=\qty{1.0}{atm}$.
\end{enumerate}
\end{corrige}

\begin{corrige}[fraction molaire]
$\chi_i=\dfrac{n_i}{n_{\text{tot}}}$ avec $n_{\text{tot}}=n_A+n_B$.
\begin{enumerate}[label=\alph*)]
  \item $n_{\text{tot}}=5$ : $\chi_A=\dfrac{2}{5}=0{,}40$, $\chi_B=\dfrac{3}{5}=0{,}60$.
  \item $\chi_A=\dfrac{0{,}25}{1{,}0}=0{,}25$.
  \item $\chi_A+\chi_B=0{,}40+0{,}60=1$ \checkmark.
\end{enumerate}
\end{corrige}

\begin{corrige}[pression partielle à partir de la fraction molaire]
$p_A=\chi_A\times p_{\text{tot}}$.
\begin{enumerate}[label=\alph*)]
  \item $p_A=0{,}4\times2{,}0=\qty{0.8}{atm}$.
  \item $p(\ce{O2})=0{,}20\times1{,}0=\qty{0.20}{atm}$.
  \item $p_B=0{,}75\times4{,}0=\qty{3.0}{atm}$.
\end{enumerate}
\end{corrige}

\begin{corrige}[fraction volumique = fraction molaire]
Pour des gaz parfaits, la fraction volumique \emph{est} la fraction molaire.
\begin{enumerate}[label=\alph*)]
  \item $\chi(\ce{O2})=0{,}20$.
  \item $\chi(\ce{N2})=0{,}80$.
  \item $p(\ce{N2})=\chi(\ce{N2})\times p_{\text{tot}}=0{,}80\times1{,}0=\qty{0.80}{atm}$.
\end{enumerate}
\end{corrige}

\end{document}
